\section{公开卖方材料：可用的是盈利分布，不可用的是当前目标价}

本案例归档了四份原始券商PDF：开源、东吴、中银和辉立。前三份均在近90天内，提供2026E/2027E收入、归母或EPS及评级；辉立的CNY147目标价基于2025E、参考价为2025年12月31日，属于历史线索。另有十二条保留券商、分析师、日期、评级和EPS字段的公开一致预期快照。

卖方材料的共同方向是看好座舱、智驾、全球化和产品创新，但研究报告不能把“买入评级”本身转成我们自己的目标价。近90天的原始PDF没有披露当前点目标价或可审计的估值方法；快照也缺目标价和方法。因此Street价格锚严格为零，卖方EPS只用作外部预期区间。

\begin{exhibitbox}[公开盈利预期与来源质量]
\small
\begin{tabularx}{\textwidth}{L{2.4cm}Y Y Y}
\toprule
来源 & 2026E EPS & 2027E EPS & 使用方式 \\
\midrule
近90天原始PDF：开源/东吴/中银 & CNY4.69--4.96 & CNY5.72--5.97 & 可作盈利预测交叉检查；无当前目标价权重 \\
辉立原始PDF（2026-01） & CNY5.23 & CNY6.63 & CNY147为基于2025E的历史目标，零权重 \\
公开结构化快照12家 & CNY4.41--5.20 & CNY5.12--6.23 & 可观察预期离散度；无收入/目标/方法字段 \\
本院基准 & CNY4.69 & 不作为远期目标输入 & 保守合并情景，等待中报验证 \\
\bottomrule
\end{tabularx}
\sourcenote{券商PDF归档、公开一致预期快照；所有币种为CNY/股。}
\end{exhibitbox}

公开快照近三个月显示八家买入、五家增持，但这一计数不是目标价共识。2026E EPS CNY4.41--5.20的区间也说明盈利预期并不完全一致。我们选择CNY4.69的基准EPS，等于原始报告区间下沿附近，并让公开中位附近预测成为上行情景的验证参照，而非默认事实。

\section{市场隐含预期与二级市场边界}

2026年7月23日盘中CNY83.48相对公开CNY4.88 EPS约为17.1倍P/E，相对本院CNY4.69基准EPS约为17.8倍。这个水平意味着市场已经为持续成长和系统能力支付价格，却尚未处于需要极端盈利爆发才合理的区间。因而本院不会因内在目标略低/略高于市场就给出机械卖出或买入，而是以“市场支持型观察”处理。

可用市场包没有可审计的250日复权序列、相对指数表现、北向/融资持仓、龙虎榜席位或严格自由流通资料。日内高低和成交只足以说明当时的价格状态，不能被称为长期支撑、阻力、趋势或资金拥挤。因此本章的二级市场结论是有限的：价格是估值锚，非技术形态结论。

\begin{keyinsight}[情绪溢价的正确处置]
若市场价格在业绩、毛利和现金验证前显著上行，须把差额标为情绪溢价，并检查其是否在破裂；若中报验证收入和利润质量，才可重新评估市场锚和倍数。没有历史资金数据时，不用“机构加仓”“北向流入”或“趋势反转”填补结论。
\end{keyinsight}

\section{与本院估值的差异}

本院CNY86.4基础业务价值并非券商目标价的替代。它由基本面CNY86.8（80\%）和市场锚CNY85.0（20\%）组合，Street目标权重为0\%。在此基础上，本院另以CNY2.49/股的Physical AI风险调整期权形成CNY88.9合并参考值；该期权同样不是券商目标价或一致预期。这一设计承认可观察的市场价格，又不把缺失的外部价格锚伪造成共识。若日后取得带日期、方法和利润假设的原始券商目标价，需重新评估其可比性，而不是追溯性地把历史值加到今天的模型。
